\documentclass[11pt]{article}
\usepackage[margin=1in]{geometry}
\usepackage[T1]{fontenc}
\usepackage{lmodern}
\usepackage{tabularx}
\usepackage{booktabs}
\usepackage{enumitem}
\setlist{nosep}

\title{ES4 Final Project Proposal: Handwritten Digit Recognizer}
\author{AJ Romeo, Peter Wiehler, Rhys de Haan, Charles Obeng-Ansong, Onyx Ward}
\date{\today}

\begin{document}
\maketitle

\section{Project description}

We will build a handwritten digit recognizer that runs entirely on our UPduino. A user draws a digit from 0 to 9 on a VGA monitor using a joystick to move a cursor and a pushbutton to paint. When they press a second ``classify'' button, a small neural network implemented in FPGA logic reads the drawing and shows its guess on the same screen. A third button clears the canvas so they can try again.

The network will be trained in Python on the MNIST dataset, then quantized to 8-bit integers (so it can fit in the FPGA RAM) and loaded onto the FPGA. All inference happens on the FPGA. The drawing is stored as a $28\times28$ grid, downsampled to $14\times14$ (196 inputs to the NN), and fed through a two-layer network ($196 \rightarrow 16 \rightarrow 10$). That is about 3,300 weights (26.4kbit), which fits in the chip's block RAM. We expect fair accuracy on MNIST test digits, with some loss on joystick-drawn digits that we will try to reduce by making strokes thick and centered.

\medskip
\noindent\textbf{Inputs:} arcade joystick, pushbuttons (paint, classify, clear).

\noindent\textbf{Outputs:} VGA monitor, driven through VGA adapter.

\medskip
\noindent\textbf{Screen layout:}
\begin{itemize}
  \item Left side: a large $28\times28$ drawing canvas, each cell drawn as a $14\times14$ pixel block ($392\times392$ pixels), with a highlighted square showing the cursor. The exact pixel values can change as long as we keep the $14\times14$ input to the network.
  \item Right side: the predicted digit drawn large, and a row of ten small bars showing the network's score for each class.
  \item Bottom: a short status line (``DRAW'', ``THINKING'', ``RESULT'').
\end{itemize}

\section{Current unknowns and point people}

\begin{tabularx}{\textwidth}{X l}
\toprule
\textbf{Unknown} & \textbf{Person} \\
\midrule
How to generate VGA timing with the PLL and draw the canvas, cursor, predicted digit and score bars on screen & Onyx \\
How to read the arcade joystick and buttons cleanly and turn them into cursor movement and painting & Charles \\
How to store the canvas in memory, downsample it, and coordinate draw, classify and clear with a top-level state machine & Rhys \\
How to build the multiply-accumulate datapath that runs both network layers using the UPduino's DSP block, including ReLU and picking the highest score & AJ \\
How to train and quantize the network in Python and export the weights into memory initialization files the FPGA toolchain can load & Peter \\
\bottomrule
\end{tabularx}

\section{Hardware needed}

One arcade joystick, three pushbuttons, a VGA adapter, an UPduino, a VGA monitor.

\end{document}
